% Standalone addition for Part A; no amenability or Ramsey claim is made here.
\begin{proposition}[Proved: an ordered vector-space parity obstruction]
\label{A:ordered-vector-parity}
Let $\mathcal V^{<}_2$ have as objects the nonzero points of finite
positive-dimensional $\mathbb F_2$-vector spaces, with an arbitrary
linear order, and as arrows the restrictions of injective linear maps
which preserve that order. This category is not exact over
$\mathbb F_2$ or over $\mathbb Z$. In fact the five embeddings
specified below have no affine equalizer over $\mathbb F_2$ at any
target.
\end{proposition}
\begin{proof}
Let $P$ be the punctured two-dimensional vector space, whose three
points are denoted $x_1<x_2<x_3$. Label the nonzero vectors of
$\mathbb F_2^3$ by $1,\ldots,7$ in binary, so addition is bitwise
exclusive-or, and give the resulting object $Q$ the reference order
\[
 1<2<3<4<6<5<7.
\]
Define $a_j:P\to Q$ by their increasing image triples:
\[
\begin{array}{c|ccccc}
 j&1&2&3&4&5\\\hline
 (a_jx_1,a_jx_2,a_jx_3)&(1,2,3)&(1,6,7)&(2,5,7)&(3,4,7)&(3,6,5).
\end{array}
\]
Each triple is the set of nonzero vectors of a two-dimensional
subspace. Since every permutation of the three nonzero vectors of
$\mathbb F_2^2$ is linear, these are arrows of the specified category.

At any target $D$, choose an auxiliary ordered vector-space basis
and let $\prec$ be its natural binary order: compare coefficient
vectors at the largest index where they differ, with $0<1$.
This auxiliary order is unrelated to the given reference order of
$D$. For a nonzero vector $v$, write $h(v)$ for its largest nonzero
basis index and let $\epsilon_v$ be the linear functional extracting
the coefficient at that index. For distinct vectors $x,y$,
\[
 \mathbf1_{x\prec y}=\epsilon_{x+y}(y)\quad\text{in }\mathbb F_2.
\]
For nonzero distinct $u,v$, a second useful identity is
\[
 \mathbf1_{u+v\text{ lies between }u\text{ and }v\text{ in }\prec}
   =\epsilon_u(v)+\epsilon_v(u).
 \tag{*}
\]
Indeed, if $h(u)=h(v)$, the sum $u+v$ precedes both and the right
side is $1+1=0$. If $h(u)>h(v)$, then $v$ precedes both $u$ and
$u+v$, and the sum lies between them precisely when
$\epsilon_v(u)=1$; also $\epsilon_u(v)=0$. The remaining case is
symmetric. These arguments work in any finite dimension.

For any linear embedding $g:Q\to D$, abbreviate $g(i)$ to $i$ in
the following calculation. The identities
$1+2=5+6=3$ and $1+6=2+5=7$ imply
\[
\begin{split}
 &\mathbf1_{2\prec1}+\mathbf1_{1\prec6}
       +\mathbf1_{2\prec5}+\mathbf1_{5\prec6}\\
 &\qquad=\epsilon_3(1)+\epsilon_7(6)
       +\epsilon_7(5)+\epsilon_3(6)
  =\epsilon_3(7)+\epsilon_7(3).
\end{split}
\]
By (*) the last expression is the indicator that $4=3+7$ lies
between $3$ and $7$. Thus the sum of these five indicators is zero
for every $g$, regardless of the reference order of $D$.

Define five functions on $\mathcal V^{<}_2(P,D)$ by
\[
\begin{array}{ll}
 f_1(e)=\mathbf1_{e(x_2)\prec e(x_1)},&
 f_2(e)=f_3(e)=\mathbf1_{e(x_1)\prec e(x_2)},\\
 f_4(e)=\mathbf1_{e(x_2)\text{ between }e(x_1),e(x_3)},&
 f_5(e)=\mathbf1_{e(x_3)\prec e(x_2)}.
\end{array}
\]
The preceding calculation says $\sum_j f_j(ga_j)=0$ for every
$g:Q\to D$. On the other hand $\sum_j f_j(e)=1$ for every
$e:P\to D$. To check the latter identity, the two equal functions
$f_2,f_3$ cancel, and the median indicator is
$\mathbf1_{e(x_1)\prec e(x_2)}+
\mathbf1_{e(x_3)\prec e(x_2)}$ modulo two.

If an affine $t\in\mathbb F_2[\mathcal V^{<}_2(Q,D)]$ equalized
the five $a_j$, write their common composite as $m$. Extend the
$f_j$ linearly. Then
\[
 0=\sum_j f_j(ta_j)=\sum_j f_j(m)
    =\operatorname{aug}(m)=\operatorname{aug}(t)=1,
\]
a contradiction. The affine criterion gives nonexactness over
$\mathbb F_2$; integral exactness would imply exactness over
$\mathbb F_2$ and is therefore impossible as well.

The attached exact finite computation enumerates all $168$ natural
orders of the seven-point space and verifies the identity. It also
checks the associated marginal matrix and all $30$ isomorphism types
of reference orders. The coordinate argument above, rather than a
bounded search over targets, proves nonexistence at every target.
\end{proof}

\paragraph{Scope of this result.}
Combined with the adjacent structural and amenability theorems, this
proposition refutes A.3(ii) and the total-category implication in
A.3(v). It does not settle rational exactness of this category or
A.3(iii); no Ramsey expansion is being assumed.
